\section{Endpoint patterns and certified overlap cutoffs}
\label{sec:endpoint-clipping}

\subsection{Extending the existing compressed-word path}
The compressed-certificates delivery adds a useful continuation of
Section~\ref{sec:lcs-endpoint}: short endpoint patterns can have regularly
spaced occurrences even when neither endpoint letter does, and a progression
of candidate overlap lengths can have a proper initial segment of successes.
The maintained implementation keeps the earlier uniform, short-period,
sparse-letter and two-largest-overlap shortcuts first. Only their unresolved
queries of length at least 128 reach the new continuation in
\path{fastunknot/compressed_endpoint.py}. A failed structural certificate
still invokes the complete occurrence-table matcher.

For words $x,y$ represented by binary straight-line programs, put
\[
 \mathcal O(x,y)=\{k:1\le k\le\min(|x|,|y|),\quad
     x[|x|-k:|x|]=y[0:k]\}.
\]
The output must describe \emph{all} these lengths by disjoint arithmetic
progressions, not merely certify one good overlap. Both operands may first
be cropped to length $L=\min(|x|,|y|)$, retaining the suffix of $x$ and
prefix of $y$. This does not change $\mathcal O$.

\subsection{Summarizing every occurrence of a short pattern}
Fix an explicit pattern $a$ of length $q$. For all its occurrence starts in
a grammar word, retain the count, first and last start, and minimum and
maximum successive gap. Empty and singleton sets have separate exact
conventions. At a concatenation $uv$, occurrences fall into three disjoint,
ordered groups: wholly in $u$, crossing the cut, and wholly in $v$.
The middle group uses at most $q-1$ possible shares from the suffix of $u$;
retaining only the first and last $q-1$ letters of each grammar word suffices
to check every such occurrence. The last group is translated by $|u|$.

Concatenating two nonempty occurrence lists adds their counts and introduces
exactly one new gap between the first list's last start and the second list's
first start. Updating both gap extrema therefore gives an exact summary by
induction. A set with at least two positions is an arithmetic progression
precisely when its two gap extrema agree. Counts and gaps are arbitrary-size
integers, not saturated counters or numerical iteration bounds.

For $g$ reachable rules, one pattern costs $O(q^2g+g\log g)$ elementary
operations and $O(qg)$ stored symbols and integer fields. The sorting term
comes from the existing reachable-rule ordering. Coordinate and label bit
lengths must still be charged in a bit model. The production trial lengths
are $q\in\{1,2,4,8\}$, in both endpoint directions; a bounded pattern prefix
is a proposal, never evidence that the rest of a word is periodic.

\subsection{The period makes candidate rank monotone}
Let $a$ be the length-$q$ suffix of $x$. Every overlap of length at least
$q$ ends at an occurrence of $a$ in $y$. Suppose these necessary candidate
lengths form
\[
 K=\{b,b+d,\ldots,b+(r-1)d\},\qquad M=\max K,
\]
with $r>1$. Check by exact compressed equality that $P=y[0:M]$ has period
$d$, meaning $P[0:M-d]=P[d:M]$. For every $k\in K$, the difference
$M-k$ is a multiple of $d$, so
\[
 y[0:k]=P[0:k]=P[M-k:M].
\]
Thus the candidate is a true overlap exactly when the length-$k$ suffixes
of $x$ and $P$ agree. Equivalently,
\[
 \mathcal O(x,y)\cap[q,\infty)
     =K\cap[\,q,\operatorname{lcsuffix}(x,P)\,].
\]
The dual argument starts with the length-$q$ prefix of $y$, takes candidate
lengths $|x|-s$ from its starts $s$ in $x$, and checks period $d$ on
$P=x[|x|-M:|x|]$. The cutoff is then $\operatorname{lcp}(P,y)$.
The first candidate need not equal the period. Material outside the maximum
candidate window is irrelevant to the period certificate.

Consequently exact overlap equality is true on an initial segment of the
candidate ranks. The implementation first tests the largest rank. If it
passes, every candidate passes. Otherwise it tests the smallest rank; failure
there rules out all candidates. In the remaining case, probe backwards from
the maximum at distances $1,2,4,\ldots$, stopping at a success or the known
minimum, then bisect the bracket. With $b_{\rm bad}$ rejected top candidates,
this takes
\[
 O(1+\log(1+b_{\rm bad}))\le O(\log(r+1))
\]
exact candidate comparisons. All-pass and all-fail cases use one and at most
two comparisons, respectively. Exactly one rejected top candidate needs three
comparisons, even when both the period and the candidate count have hundreds
of bits. Computing a full LCP on a much larger numerical range is unnecessary.

An empty occurrence set excludes all lengths at least $q$. A singleton
candidate is checked directly, without a period assumption. Every shorter
length $1,\ldots,q-1$ is checked separately. A successful trial therefore
returns the complete overlap set: at most one progression of long overlaps
and at most $q-1$ short singletons. A nonprogression or failed exact period
check gives no conclusion and tries another bounded anchor, then falls back.

\subsection{Costs, integration and examples}
If $H$ bounds the input grammar height, $c$ is the rank-comparison count and
$E(G)$ bounds exact equality at the actual intermediate grammar size, a
conservative successful-trial cost is
\[
 O(q^2g+g\log g+H(q+c))+O(q+c+1)E(G),
 \qquad G=O(g+H(q+c))
\]
for a fresh arena containing the input and query artifacts. Earlier failed
trials have their own bounded preparation and equality costs. This is an
operation bound, not a unit-cost assumption for binary positions or letter
labels. In particular, a middle cutoff can still require many slices and
exhaust the node allowance. No general linear-time claim is made.

All preparation and exact checks use the caller's existing word-arena budget.
Local resource exceptions and global cancellation propagate. Neither becomes
an empty overlap answer; incomplete queries are not inserted in the overlap
cache. The old endpoint counters are preserved, while new preparation,
period, cutoff and allocation counters have the \path{lcs_anchor_} prefix.
The existing longest-common-substring critical-extension calculation receives
the same complete progression interface. Forty separately checked clipped
families verify that its finite critical candidates still attain the maximum
over every explicit overlap.

For $n,m\ge2$, the words
\[
 ((21)^n1)^m,\qquad ((12)^n11)^m
\]
have dense single-letter endpoints. A suffix bigram in the first case gives
all positive multiples of $2n+1$ through $m(2n+1)$. A four-letter anchor in
the second gives all positive multiples of $2n+2$ through $m(2n+2)$, together
with the short overlap one. Their grammars have $O(\log n+\log m)$ rules.
For $P_n=(12)^n3$, the pair
\[
 x=4^{2n+1}P_n^{m-1},\qquad y=P_n^m
\]
has exactly one rejected top candidate, so the three-comparison cutoff applies.
The older all-success marker family remains on the original two-largest path.
A damaged word whose relevant windows have no certified period still needs
the general matcher; the extension is not a universal pattern shortcut.

\subsection{Validation and comparison with the current baseline}
The 12 new test methods cover exact signed occurrence summaries, exhaustive
short-word overlaps, 360 random parse pairs, eight-letter anchors, a depth-4000
grammar, 500-bit periods and counts, interior and near-end cutoffs, allocation
and work limits, and cancellation without partial cache entries. The four
previous endpoint-progression tests still pass, including a trap proving that
their successful old shortcut never invokes the new module. Existing general
fallback tests explicitly disable the additional optional shortcut so that
they continue exercising the table path. All 768 maintained tests pass in
167.676 seconds with Regina installed.

The benchmark fixes the old matcher at \code{8a9c08bcf} inside otherwise
common current word, relator and recognition code. It runs old, identical
A/A control and maintained clipping arms, with one excluded warmup followed
by five measured rounds in shuffled order. Source hashes are frozen before
measurement and checked afterward. Each compressed operation uses fresh
state, two million shared work units and at most 100,000 grammar nodes.
Grammar construction is included in timings; output containment and
cardinality checks are outside the overlap-query clock. Numerical word
lengths are never used to enumerate the large expected answers.

Eight overlap families use parameters $n=2^b$ at $b=8,64,500$, covering
new dense endpoints and cutoff classes, successful old shortcuts, and
uncertifiable periods. Whole phased LCS and cyclic relator searches use the
dense bigram family at $b=8,32$; the relator scope includes failed whole-donor
search, the complete cyclic overlap query and checked application. A pure-power
quotient control retains the established path. These are compressed algebraic
inputs, not diagrams with $2^b$ crossings.

\begin{center}\small
\begin{tabular}{@{}lrrrrr@{}}
\toprule
Family & Bits & Before & A/A & Clipping & Ratio \\
\midrule
Marker control & 8 & 0.821 & 0.862 & 0.870 & 0.94 \\
Marker control & 64 & 6.120 & 6.187 & 6.046 & 1.00 \\
Marker control & 500 & 51.187 & 50.593 & 52.121 & 1.00 \\
Dense bigram & 8 & 307.677 & 310.654 & 1.259 & 244.48 \\
Dense bigram & 64 & limit & limit & 8.145 & -- \\
Dense bigram & 500 & limit & limit & 65.936 & -- \\
Dense fourgram & 8 & 308.246 & 313.962 & 1.530 & 201.36 \\
Dense fourgram & 64 & limit & limit & 10.193 & -- \\
Dense fourgram & 500 & limit & limit & 85.070 & -- \\
One rejected & 8 & 329.133 & 344.696 & 1.198 & 279.34 \\
One rejected & 64 & limit & limit & 7.965 & -- \\
One rejected & 500 & limit & limit & 62.965 & -- \\
Half rejected & 8 & 308.250 & 303.671 & 2.076 & 150.16 \\
Half rejected & 64 & limit & limit & 62.613 & -- \\
Half rejected & 500 & limit & limit & limit & -- \\
Period rejected & 8 & 384.168 & 378.957 & 378.362 & 1.02 \\
Period rejected & 64 & limit & limit & limit & -- \\
Period rejected & 500 & limit & limit & limit & -- \\
Short-period control & 8 & 0.090 & 0.091 & 0.099 & 0.91 \\
Short-period control & 64 & 0.353 & 0.331 & 0.330 & 1.07 \\
Short-period control & 500 & 2.243 & 2.281 & 2.243 & 1.00 \\
Sparse control & 8 & 0.798 & 0.800 & 0.796 & 1.00 \\
Sparse control & 64 & 1.139 & 1.122 & 1.126 & 1.01 \\
Sparse control & 500 & 3.737 & 3.701 & 3.753 & 0.99 \\
\bottomrule
\end{tabular}
\end{center}
Median milliseconds; ratios are medians of complete paired before/clipping times. Censored runs are shown as limits and excluded from ratios.

\begin{center}\small
\begin{tabular}{@{}lrrrrr@{}}
\toprule
Family & Bits & Before & A/A & Clipping & Ratio \\
\midrule
Phased LCS & 8 & 940.259 & 934.837 & 59.307 & 15.88 \\
Phased LCS & 32 & limit & limit & 1053.355 & -- \\
Phased relator & 8 & limit & limit & 155.734 & -- \\
Phased relator & 32 & limit & limit & limit & -- \\
Quotient control & 8 & 0.344 & 0.340 & 0.335 & 1.02 \\
Quotient control & 32 & 0.963 & 0.962 & 0.962 & 1.00 \\
\bottomrule
\end{tabular}
\end{center}
Median milliseconds; ratios are medians of complete paired before/clipping times. Censored runs are shown as limits and excluded from ratios.


The complete-query comparison retains all 15 established hard-unknot inputs,
all existing filters, compressed group search and certificate replay. The
new two-meridian stage is disabled so that it does not bypass the matcher
being tested. Every query starts from a freshly validated PD. Group work is
capped at 20 million units and 15 seconds, with an 18-second global allowance
and 50,000 Khovanov objects. Certificate digests and complete search counters
are retained outside the timed region. Timeout records never enter
completed-time medians or paired speedup denominators.

\begin{center}\small
\begin{tabular}{@{}lrrrr@{}}
\toprule
Input & Before & A/A & Clipping & Ratio \\
\midrule
Survivor 00 & 17.824 & 17.633 & 17.360 & 1.00 \\
Survivor 01 & 14.852 & 14.786 & 14.795 & 1.00 \\
Survivor 02 & 16.129 & 16.147 & 16.122 & 1.00 \\
Survivor 03 & 11.309 & 11.290 & 11.266 & 1.00 \\
Mirror 03 & 24.843 & 24.946 & 25.149 & 0.99 \\
Survivor 04 & 9.927 & 10.000 & 9.954 & 1.00 \\
Survivor 05 & 19.091 & 19.155 & 19.104 & 1.00 \\
Survivor 06 & 7.834 & 7.662 & 7.626 & 1.00 \\
Survivor 07 & 19.815 & 19.931 & 19.789 & 1.00 \\
Survivor 08 & 7.694 & 7.674 & 7.687 & 1.00 \\
Mirror 08 & 20.825 & 20.633 & 20.695 & 1.02 \\
Survivor 09 & 7.600 & 7.528 & 7.502 & 1.00 \\
Survivor 10 & 6.518 & 6.513 & 6.650 & 0.97 \\
Survivor 11 & 7.210 & 7.122 & 7.142 & 1.01 \\
Gordian & 1968.756 & 1988.703 & 1972.944 & 1.00 \\
\bottomrule
\end{tabular}
\end{center}
Median milliseconds; ratios are medians of complete paired before/clipping times. Censored runs are shown as limits and excluded from ratios.


All raw samples and source hashes are retained in:
\begin{center}\small
\path{fast/results/endpoint_clipping_20261008.json}
\end{center}
The record contains 675 measured calls: 360 complete-overlap attempts, 90
whole algebraic operations and 225 whole-knot queries, plus 135 excluded
warmups. Of the measured calls, 525 complete and 150 exhaust an allowance.
All completed answers agree with their structural oracle or prior knot
label; every completed knot arm produces the same certificate digest as its
paired competitors.

At the eight-bit parameter, dense bigram and fourgram overlap queries improve
by paired factors $244.48$ and $201.36$, respectively. The one-rejected family
improves by $279.34$ and the half-rejected family by $150.16$. At 500 bits,
dense bigram, dense fourgram and one-rejected queries complete in median
65.94, 85.07 and 62.96 milliseconds; both old arms exhaust the work cap.
Those capped comparisons demonstrate increased capacity, not numerical
speedup ratios. The one-rejected query uses exactly three rank comparisons
at all three parameter sizes. The established all-success marker, short-period
and sparse controls never invoke the new continuation.

The finite comparison also exhibits its limits. At 64 bits the half-rejected
family completes with 128 rank comparisons; at 500 bits it exhausts the
100,000-node allowance after 72 comparisons, before finding its cutoff.
The failed-period family remains on the complete fallback and still exhausts
the work cap at 64 and 500 bits. These are unresolved matching computations,
not certificates that the overlap set is empty.

The whole phased LCS improves from 940.26 to 59.31 milliseconds at eight bits,
a paired factor $15.88$, and completes at 32 bits where the old arms exhaust
the cap. The full phased cyclic-relator query and checked application complete
in 155.73 milliseconds at eight bits where both old arms are capped. At 32
bits that operation remains capped even with eight successful anchor hits:
other matching and extension work still dominates the shared allowance.
The quotient control remains on its old path.

All 225 measured whole-knot queries complete, but their new anchor activation
counts are zero. Their before/clipping ratios range from $0.967$ to $1.018$,
while their A/A ratios range from $0.993$ to $1.025$. There is consequently no
measured recognition speedup on this corpus. For completed primitive pairs,
A/A ratios range from $0.961$ to $1.067$; small changes on unaffected controls
are retained as timing variation. The supported conclusion is broader exact
compressed-word capacity and faster selected complete algebraic queries.
Neither a general relator-search move bound nor a quasi-polynomial knot
recognizer follows from these results.
